\section{Contamination Checks for Loop Engineering Evaluation}\label{app:contamination}

\lhb{} derives its loop engineering evaluation tasks from public development artifacts, including PR Sequences, Course Labs, and Research Evolutions. This source choice creates potential overlap with modern LLM pretraining data. We therefore separate source recognition risk from the loop trace metrics used to evaluate loop execution.

\paragraph{Input Side Anonymization.}
During Task Instrumentation (Section~\ref{sec:description-refine}), every unit requirement is rewritten to remove identifiable tokens that could let a model recognize the source from the prompt: repository names, course identifiers (e.g., MIT 6.824, Stanford CS144), paper titles, and author attributions. For PR Sequences, the same rewriting policy is applied at the segment (SEG) level, where original commit messages and PR titles often include the feature name, issue tracker ID, or merging author. After this step, the task prompt provided to the evaluated loop contains the acceptance specification, public interfaces, and tests. Source identifying evidence that would support retrieval by name is removed before evaluation.

\paragraph{Memorization and Loop Trace Metrics.}

Anonymization addresses prompt level identification, but training set memorization may remain in model weights. A model that has internalized a reference implementation, for example, may reproduce local algorithmic fragments regardless of prompt wording. Two observations limit the extent to which this residual risk explains the loop trace metrics recorded by \lhb{}.

The first observation is empirical. The strongest configuration in our evaluation, Claude Code with Opus-4.7, uses the native Anthropic model and loop infrastructure under the most permissive context budget and reaches a Resolve Rate of $25.00\%$ (Table~\ref{tab:rq1}). If weight level memorization were the principal driver, frontier model and loop configurations would be expected to approach higher resolution on canonical Course Labs and Research Evolutions tasks, whose source material is publicly accessible. This result provides limited support for memorization as the main explanation.

The second observation is structural. \lhb{} treats recall of a single implementation fragment as insufficient evidence for loop engineering evaluation. A resolved task entails prerequisite consistency across many edits, retained obligations as the codebase evolves, preserved state across context renewal, and test feedback routed along the dependency DAG. These loop conditions extend beyond pretraining exposure to a solution for a single local unit. Memorization may supply fragments, but \lhb{} records whether those fragments remain consistent under live regression obligations, where current model and loop configurations exhibit constrained state continuity and obligation retention.

\paragraph{Model cutoff robustness slice.}

A direct slice by source ingest date and model cutoff can partition tasks by whether their source material predates each model's training cutoff and report Resolve Rate within each partition. The current analysis reports the aggregate constraint that the strongest model and loop configuration reaches $25\%$ Resolve Rate, providing limited support for memorization driven inflation as the principal explanation.
